% Generated draft: requires mathematical and source review.
\documentclass[11pt]{article}
\usepackage[margin=25mm]{geometry}
\usepackage{fontspec}
\setmainfont{Latin Modern Roman}
\usepackage{amsmath,amssymb,mathtools,unicode-math}
\setmathfont{Latin Modern Math}
\usepackage{enumitem,needspace,xurl,hyperref}
\setlength{\parindent}{0pt}
\setlength{\parskip}{5pt}
\setlength{\emergencystretch}{3em}
\newcommand{\N}{\mathbb{N}}
\newcommand{\Z}{\mathbb{Z}}
\newcommand{\Q}{\mathbb{Q}}
\newcommand{\R}{\mathbb{R}}
\newcommand{\C}{\mathbb{C}}
\newcommand{\F}{\mathbb{F}}
\newcommand{\A}{\mathbb{A}}
\newcommand{\PP}{\mathbb{P}}
\newcommand{\E}{\mathbb{E}}
\newcommand{\eps}{\varepsilon}
\newcommand{\dd}{\,\mathrm{d}}
\begin{document}
\section*{UnsolvedMath ID: 30002191}
\textbf{Source problem code:} catalogue entry
\subsection{Definitions and mathematical statement}
The supplied record asks for an identity relating nested punctual Hilbert schemes of a planar curve singularity \(C\) at \((0,0)\) to the triply graded HOMFLY homology of its algebraic link. Write \(\mu\) for its Milnor number and \(a,q,t\) for formal grading variables.

The base field, reducedness, and isolatedness hypotheses are not explicitly specified. Under the usual interpretation as a reduced complex plane curve germ \(C=\{f=0\}\subset(\C^2,0)\) with an isolated singularity, its algebraic link is
\[
L_{C,(0,0)}=C\cap S^3_{\varepsilon},
\]
where \(S^3_{\varepsilon}\) is a sufficiently small sphere centered at the origin, and its Milnor number is
\[
\mu=\dim_{\C}\frac{\C\{x,y\}}{(\partial f/\partial x,\partial f/\partial y)}.
\]
These hypotheses describe a conventional interpretation; the supplied record does not explicitly confirm them or specify the orientation convention for the link.

The notation \(C_{(0,0)}^{[\ell,\ell+m]}\) refers to a nested punctual Hilbert scheme. For ordinary nesting, its points would be pairs \(Z\subseteq Z'\) of finite subschemes of \(C\), both supported at the origin, with respective lengths \(\ell\) and \(\ell+m\). Equivalently, their ideals satisfy \(I_{Z'}\subseteq I_Z\) and
\[
\dim_{\C}(\mathcal O_{C,0}/I_Z)=\ell,
\qquad
\dim_{\C}(\mathcal O_{C,0}/I_{Z'})=\ell+m.
\]
The record does not define the intended incidence conditions, so this ordinary nesting description does not establish that it is the scheme intended by the formula. Any additional condition involving the maximal ideal must be specified from the source. The notation \(P_t^{\mathrm{vir}}\) denotes the source's virtual Poincar\'e polynomial, and \(P(L_{C,(0,0)})\) denotes the source's triply graded HOMFLY homology series; their precise conventions are not supplied.

\textbf{Proposed identity in the supplied statement and clean statement.} For every planar curve singularity in the intended, but incompletely specified, class, the proposed identity is
\[
\left(\frac aq\right)^{\mu-1}
\sum_{\ell,m}q^{2\ell}a^{2m}t^{m^2}
P_t^{\mathrm{vir}}\!\left(C_{(0,0)}^{[\ell,\ell+m]}\right)
=P(L_{C,(0,0)}).
\]
The summation ranges are not stated. Nonnegative integers \(\ell,m\) would be consistent with the length interpretation above, but that range requires confirmation.

\textbf{Conflicting supplied formulation.} The original-statement field instead gives
\[
\sum_{\ell,m}\left(\frac aq\right)^{\mu-1}
q^{2\ell}a^{2m}t^m
P^{\mathrm{vir}}\!\left(C_{(0,0)}^{[\ell,\ell+m]}\right)
=P(L_{C,(0,0)}).
\]
In particular, the explicit factor is \(t^m\) here and \(t^{m^2}\) in the other formulation; the polynomial notation also differs. The supplied evidence does not explain a change of convention that reconciles these expressions. Thus the intended identity cannot be specified uniquely without resolving this discrepancy and fixing the incidence, polynomial, homology, and grading conventions. No equality or mathematical resolution is asserted by this draft.
\subsection{Short English statement}
The proposed conjecture identifies a generating series of virtual Poincar\'e polynomials of nested punctual Hilbert schemes with the triply graded HOMFLY homology series of the singularity's algebraic link. The supplied formulas disagree between \(t^{m^2}\) and \(t^m\), and several hypotheses and normalization conventions remain unspecified. The supplied status is partially solved; this is dated curation, not an independently verified conclusion.
\subsection{Sources}
\begin{itemize}
\item Supplied local UnsolvedMath registry record, permanent ID \(30002191\), source problem code \texttt{catalogue entry}, title ``Homological Hilbert-Scheme Conjecture for Algebraic Links,'' in the Oberwolfach Reports Open Problems collection. The record gives proposed year 2012 and names no proposer.
\item Supplied report citation: \emph{Singularities (2013)}, Oberwolfach Reports, DOI \texttt{10.4171/owr/2012/46}, pp.\ 2799--2866, \url{https://doi.org/10.4171/owr/2012/46}. This is also the supplied literature-source link labeled ``original OWR report.'' The report title's year 2013 and the record's proposed year 2012 are retained as supplied.
\item Supplied literature source: Oblomkov--Rasmussen--Shende, ``HOMFLY homology conjecture,'' \url{https://arxiv.org/abs/1201.2115}.
\item Supplied literature source: ``Khovanov--Rozansky homology for Coxeter knots,'' \url{https://arxiv.org/abs/2407.18123}.
\item Literature assessment dated exactly \texttt{2026-08-21}: the supplied curation labels the problem partially solved and states that the homological Hilbert-scheme formula remains open in full generality, with substantial families and specializations proved, including torus-knot cases and related geometric realizations. This is supplied dated curation, not independently verified literature assessment; the record expressly requests independent verification of time-sensitive status and cited claims.
\item The registry labels the statement \texttt{corrected\_verified} and reports a check against the official Oberwolfach report. That is a supplied verification claim, not a check performed for this draft; the conflict with the original-statement field remains unresolved. No source was visited in preparing this statement-only draft.
\end{itemize}
\end{document}
